\chapter{Regenerating Theoretical Bounds}
\label{ch:regenerating_bounds}

Throughout this text, we have presented numerous theoretical bounds, visualizing them with \texttt{pgfplots}. However, these visualizations are most instructive when grounded in actual numerical computation rather than idealized or sketched curves. In this chapter, we provide the corresponding Python scripts to regenerate the numerical data for these bounds. By executing these snippets---designed to run with standard, dependency-light numerical libraries such as \texttt{numpy}, \texttt{scipy}, and \texttt{torch}---you can export precise coordinate files (e.g., CSVs or text files) that directly feed into the \texttt{pgfplots} figures seen earlier.

\section{Rate-Distortion Bounds for Gaussian Sources}

Our discussion of the rate-distortion theorem in Section~\ref{sec:rate_distortion_theory} culminated in the classic bound for a memoryless Gaussian source, $R(D) = \frac{1}{2} \log_2(\sigma^2 / D)$ for $D \leq \sigma^2$. To accurately reproduce the rate-distortion curve in Figure~\ref{fig:gaussian_rd_curve}, we must compute this bound over a range of distortion values and export the results. The following snippet calculates the theoretical rate for a standard Gaussian source ($\sigma^2 = 1$) across a logarithmic grid of distortions and saves the coordinates to a file that \texttt{pgfplots} can ingest via the \texttt{addplot table} directive.

\begin{pycode}
import numpy as np

def compute_gaussian_rd(sigma_sq=1.0, num_points=100):
    # D must be in (0, sigma_sq]
    D_vals = np.logspace(-3, 0, num_points)
    R_vals = 0.5 * np.log2(sigma_sq / D_vals)
    
    with open("gaussian_rd_data.csv", "w") as f:
        f.write("D,R\n")
        for d, r in zip(D_vals, R_vals):
            f.write(f"{d:.5f},{r:.5f}\n")

compute_gaussian_rd()
\end{pycode}

\section{Information Bottleneck Generalization Bounds}

In Chapter~\ref{ch:ib_generalization}, Theorem~\ref{thm:ib_gen_bound} established a generalization bound based on the Information Bottleneck (IB) principle, where the generalization gap is bounded by $\mathcal{O}(\sqrt{I(W; Z) / n})$. Here, $W$ represents the weights, $Z$ the training sample, and $n$ the dataset size. To visualize how this bound scales with sample size and mutual information complexity, the script below simulates the bound's trajectory for various hypothetical mutual information values. It uses \texttt{torch} to efficiently compute the tensor operations required for grid evaluations, producing a dataset that reproduces the surface plot of Figure~\ref{fig:ib_gen_surface}.

\begin{pycode}
import torch

def compute_ib_bounds():
    # Grid of sample sizes n and mutual information I(W; Z)
    n_vals = torch.linspace(100, 10000, 50)
    I_vals = torch.linspace(1.0, 50.0, 50)
    
    N, I = torch.meshgrid(n_vals, I_vals, indexing='ij')
    
    # Assume a constant C for the O() notation
    C = 2.0
    bound = C * torch.sqrt(I / N)
    
    with open("ib_gen_bound.csv", "w") as f:
        f.write("n,I,bound\n")
        # Flatten and save for pgfplots 3D surf plot
        for i in range(N.shape[0]):
            for j in range(N.shape[1]):
                f.write(f"{N[i,j].item():.1f},{I[i,j].item():.1f},{bound[i,j].item():.5f}\n")

compute_ib_bounds()
\end{pycode}

\section{PAC-Bayesian Complexity Bounds}

Finally, we turn to the PAC-Bayesian bound introduced in Equation~\ref{eq:pac_bayes_kl}. This bound rigorously limits the expected risk of a stochastic classifier using the Kullback-Leibler divergence between a learned posterior $Q$ and a prior $P$ over the hypothesis space. To recreate the contour plot in Figure~\ref{fig:pac_bayes_contours}, which illustrates the trade-off between the empirical risk and the KL penalty, we evaluate the standard PAC-Bayes McAllester bound: $\mathcal{R}(Q) \leq \hat{\mathcal{R}}_n(Q) + \sqrt{\frac{D_{\text{KL}}(Q \| P) + \ln(2\sqrt{n}/\delta)}{2n}}$. The script generates this bound over a grid of empirical risks and KL divergences using \texttt{scipy} for robust statistical constants and \texttt{numpy} for vectorized math.

\begin{pycode}
import numpy as np
from scipy import stats

def compute_pac_bayes_contour(n=1000, delta=0.05):
    emp_risk = np.linspace(0.0, 0.5, 40)
    kl_div = np.linspace(0.0, 10.0, 40)
    
    R_hat, KL = np.meshgrid(emp_risk, kl_div)
    
    # PAC-Bayes McAllester bound
    penalty = np.sqrt((KL + np.log(2 * np.sqrt(n) / delta)) / (2 * n))
    true_risk_bound = R_hat + penalty
    true_risk_bound = np.clip(true_risk_bound, 0, 1.0)
    
    with open("pac_bayes_contour.csv", "w") as f:
        f.write("EmpRisk,KL,TrueRiskBound\n")
        for i in range(R_hat.shape[0]):
            for j in range(R_hat.shape[1]):
                f.write(f"{R_hat[i,j]:.4f},{KL[i,j]:.4f},{true_risk_bound[i,j]:.4f}\n")

compute_pac_bayes_contour()
\end{pycode}
